\subsection{Cover execution and timing-informed removal}\label{sec:cover-study}
Two fixed-weight models, Qwen3:8B and Llama3.1:8B, each run 300 synthetic ticket fixtures per condition: 120 training, 60 validation and 120 test, balanced across six actions. The six conditions yield 3,600 episodes with terminal outcomes, including 1,440 held-out executions. Fixtures and cover draws are paired; request phrasings recur across splits.

\paragraph{Execution and observations}
LangGraph~\cite{langchain2025langgraph} lets each model choose checks, actions, objects, recovery and finish. Seven local MCP Streamable HTTP destinations use TLS~1.3/HTTP/2, a common \path{/mcp} path, and 4,096/16,384-byte request/response padding quanta. A frozen schedule places one or two decoys among early calls. Policies are no cover, proxy discard, read cover, hidden isolated execution, live writes, and visible isolated execution. Proxy discard computes on a disposable copy; hidden isolation retains a clone. Live writes deliberately change the primary fixture as a state-pressure control. Visible isolation forwards clearly marked cover responses to the planner.

The original observer receives destinations and bidirectional ciphertext counts in supplied episode/call windows. LR/RF fit each condition, with validation selecting the primary attacker; position and frozen attacks remain reported. Goal completion checks the requested action/object, prerequisites and target postcondition. Whole-state integrity additionally checks other orders, shared quota and the ledger against the intended-action reference. Each original endpoint retains its twelve-test Holm family. Intervals resample paired fixtures by action, conditional on fitted attacks and the designed workload.

\begin{table*}[t]\centering\small
\caption{Same held-out cover executions, 120 paired fixtures per model/condition. Byte is the original validation-selected size/destination attack; Last is the last-destination rule. Strip adds client-receipt timing to detect/remove cover, then uses the frozen no-cover classifier (planned secondary analysis). Goal, State and Goal+error retain the original execution endpoints; call bytes include both directions. Chance recognition is 20/120.}
\label{tab:cover-main}
\setlength{\tabcolsep}{4pt}
\begin{tabularx}{\textwidth}{@{}lYrrrrrrr@{}}\toprule
Model & Cover policy & Byte & Last & Strip & Goal & State & Goal+error & Bytes $\times$ \\\midrule
Qwen3:8B & No cover & 118 & 116 & 118 & 114 & 116 & 0 & 1.00 \\
Qwen3:8B & Proxy discard & 62 & 77 & 114 & 114 & 116 & 0 & 1.34 \\
Qwen3:8B & Read cover & 62 & 77 & 114 & 114 & 115 & 0 & 1.34 \\
Qwen3:8B & Isolated, hidden & 61 & 77 & 112 & 114 & 115 & 0 & 1.34 \\
Qwen3:8B & Live write & 62 & 75 & 111 & 102 & 7 & 95 & 1.34 \\
Qwen3:8B & Isolated, visible & 63 & 77 & 111 & 113 & 114 & 1 & 1.35 \\
\midrule
Llama3.1:8B & No cover & 113 & 79 & 113 & 82 & 69 & 13 & 1.00 \\
Llama3.1:8B & Proxy discard & 79 & 67 & 115 & 80 & 68 & 13 & 1.29 \\
Llama3.1:8B & Read cover & 72 & 67 & 112 & 84 & 71 & 14 & 1.28 \\
Llama3.1:8B & Isolated, hidden & 82 & 67 & 112 & 83 & 70 & 13 & 1.28 \\
Llama3.1:8B & Live write & 80 & 65 & 112 & 76 & 3 & 73 & 1.29 \\
Llama3.1:8B & Isolated, visible & 81 & 68 & 107 & 83 & 73 & 11 & 1.29 \\
\bottomrule\end{tabularx}
\end{table*}


\paragraph{Padded calls still support removal}
In the original view, the selected attack scores 61--63/120 for Qwen and 72--82/120 for Llama under cover. The last-destination rule reaches 75--77/120 for Qwen, exposing the schedule's uncovered tail. We next add \emph{client-receipt} call starts/ends. A per-model detector learns honest/cover labels from training traces and selects its classifier/threshold on validation. Its inputs are call gaps, durations, destination and the known first-call flag. It removes predicted decoys, then applies the validation-selected no-cover classifier. Test cover truth is used only for scoring and a privileged reference.

The detector identifies 1,807/1,820 decoys and removes none of 6,741 honest calls. Task recognition becomes 111--114/120 for Qwen and 107--115/120 for Llama (Table~\ref{tab:cover-main}, Strip). All ten planned secondary removal contrasts survive Holm correction, with adjusted $p\leq6.17\times10^{-6}$. The ten primary contrasts that simply concatenate timing with the original features do not reject and show no improvement. This distinction matters: timing can reveal \emph{which calls to discard} without helping a direct task classifier. The analysis was frozen before new test-timing evaluation; original M3 test scores were already public. Supplement S9 preserves both test families and the separately labelled post-scoring ablation.

\paragraph{The controller supplies the timing distinction}
The implementation calls its decoy dispatcher directly after a tool response. Honest calls usually follow a model decision; inserted cover skips that wait. The resulting gap helps separate the two call types despite body padding. Client receipts include local scheduling and a short counting wait; this is a measured endpoint-timing attack, with call windows and cover-labelled training examples supplied. It does not establish packet-level recovery under network jitter. It does establish why early/tail coverage alone is an insufficient next test: a revised schedule must also face timing-informed removal.

\paragraph{Business outcomes remain a separate constraint}
Proxy, read and hidden-isolated cover retain Qwen state integrity at 115--116/120; adverse live writes reduce it to 7/120. Llama falls from a 69/120 baseline to 3/120. Both live-write integrity losses survive the original correction; completion loss survives only for Qwen. Of 102 Qwen and 76 Llama completed live-write tasks, 95 and 73 have collateral state errors. These counts measure the deliberately adverse control, rather than production cover harm. Call bytes grow by about 1.28--1.35$\times$. Similar byte-view scores do not establish policy equivalence, and the new removal attack further limits their use as a privacy operating point.

Making isolated responses visible raises mean prompt tokens by 18.1\% and 18.9\%; all original visible-minus-hidden average endpoint tests have adjusted $p=1$. Paths vary between paired runs, so this control remains inconclusive (Supplement S8).
